\documentclass[xcolor=svgnames,17pt]{beamer}

\usepackage[export]{adjustbox}
\usepackage{bookmark}
\usepackage{colortbl} \arrayrulecolor[gray]{0.7}
\usepackage{microtype}
\usepackage{pgfpages}
\usepackage{rotating}
\usepackage{textcomp}
\usepackage{tabularx}
\usepackage{xspace}

\usepackage{fontspec}

\hypersetup{hidelinks,pdfpagemode=}

\urlstyle{same}

\newcommand*{\sizefont}[1]{%
    \ifcase#1\relax
    \or \tiny
    \or \scriptsize
    \or \footnotesize
    \or \small
    \or \normalsize
    \or \large
    \or \Large
    \or \LARGE
    \or \huge
    \or \Huge
    \fi}

%%

\newcommand*{\mybullet}{\tikz[baseline=-.6ex]\node[%
    draw,circle,inner sep = -0.15ex,fill]{.};\xspace}

\setbeamertemplate{footline}{
    \usebeamercolor[fg]{page number in head/foot}%
    \usebeamerfont{page number in head/foot}%
    \hspace*{1ex}\insertframenumber\,/\,\inserttotalframenumber\hfill
    github.com/andrewdotn/extending-jenkins-with-jython\ }

\newcommand*{\plainfooter}{%
    \setbeamertemplate{footline}{
        \usebeamercolor[fg]{page number in head/foot}%
        \usebeamerfont{page number in head/foot}%
        \hspace*{1ex}\insertframenumber\,/\,\inserttotalframenumber\vskip2pt}}

\makeatletter
\def\alphslide{\@alph{\intcalcAdd{1}{\intcalcSub{\thepage}{\beamer@framestartpage}}}}
\newcommand*{\plainstepfooter}{
    \setbeamertemplate{footline}{
        \usebeamercolor[fg]{page number in head/foot}%
        \usebeamerfont{page number in head/foot}%
        \hspace*{1ex}\insertframenumber\alphslide\,/\,\inserttotalframenumber\vskip2pt}}
\makeatother

\setbeamertemplate{note page}{
    \sizefont{3}
    \setlength{\parskip}{10pt}
    \insertnote
    \par}

\setbeamertemplate{navigation symbols}{}
\setbeamerfont{title}{size=\LARGE}
\setbeamerfont{frametitle}{size=\LARGE}
\setbeamerfont{framesubtitle}{size=\normalsize}

\newcommand*{\tocsection}[1]{\pdfbookmark[2]{#1}{#1}}

%%

\title{Extending Jenkins \\ with Jython}

\author{\texorpdfstring{%
    Andrew Neitsch}{Andrew Neitsch}}

\date{\small 2015-10-23}

\begin{document}

\tocsection{Title page}

\begin{frame}[plain]
\titlepage
\adjustbox{width=0.9\paperwidth,center}{\structure{\href{https://github.com/andrewdotn/extending-jenkins-with-jython}{github.com/andrewdotn/extending-jenkins-with-jython}}}
\end{frame}

\note{This is a note}

\begin{frame}{Outline}
\tableofcontents
\end{frame}

\section{Introduction}

\begin{frame}{Jenkins}

\begin{itemize}
\item Industry-standard continuous integration tool
\item Regularly and automatically builds and tests your code
\item Helps make sure your code works, and stays working
\end{itemize}

\end{frame}

\begin{frame}{Example: Testing hello}
% TODO: screenshot
\end{frame}

\begin{frame}{}
If you have tests, you’re already running them before you commit, so why
bother?
\pause
\begin{itemize}
\item Can test regularly forever, not just when there’s a change
\pause
\item Multi-configuration builds: different Python versions, library
versions, platforms, Linux, Windows, and Mac
\end{itemize}
\alert{These are really valuable things that nobody will consistently do by
hand!}
\end{frame}

\begin{frame}{Example: hello v2, using six}

\end{frame}

\begin{frame}{Advanced uses}

\begin{columns}
\column{.4\textwidth}
\fbox{%
    \includegraphics[width=\textwidth]{jenkins-the-definitive-guide.png}}
\column{.6\textwidth}
\begin{itemize}
\item Try servers
\item Track betas and release candidates
\item Automated deployments
\end{itemize}
\end{columns}

\end{frame}

\section{Test Isolation}

\begin{frame}{}
\tableofcontents[currentsection]
\end{frame}

\begin{frame}[plain]{}
\includegraphics[height=0.95\paperheight,center]{works-on-my-machine.pdf}
\end{frame}

\section{Extending Jenkins}

\begin{frame}{}
\tableofcontents[currentsection]
\end{frame}

\begin{frame}{Plugins}
Over 1000 plugins
\end{frame}

\section{Conclusions}

\begin{frame}{}
\tableofcontents[currentsection]
\end{frame}

\begin{frame}{Conclusions}
\begin{itemize}
\item Jenkins runs tests that humans don’t have the patience to run by
hand: multiple Python versions, multiple platforms, every commit, every day
\pause
\item You can help your fellow developers by making tests that will run
anywhere
\pause
\item You can extend Jenkins with Jython ...
\pause but it’s probably not worth it
\end{itemize}
\end{frame}

\begin{frame}
\adjustbox{width=0.5\paperwidth,center}{\structure{Questions?}}
\end{frame}

\begin{frame}{Call for talks}

“w/o interesting talks, there's not a ton of point in ‘meeting up’”

\pause

\begin{enumerate}
\item Pick something you find interesting
\item Talk about it
\item Include exercises for people to hack on
\end{enumerate}

\end{frame}

\begin{frame}{Suggested exercises}
\begin{itemize}
\item Install Jenkins, and get it testing \\ some code you care about
\item Isolate your tests
\item Try extending Jenkins
\end{itemize}
\end{frame}

\end{document}
